% --- Archon physics formalization source begin ---
% archon:physics
% archon:covers PhyXMiniProblems/problem_phyx_mini_0593.lean
% archon:source-report reports/phyx_mini/problem_phyx_mini_0593.source.json

\chapter{Physics problem phyx\_mini\_0593}
\label{ch:PhyXMiniProblems_problem_phyx_mini_0593}

\paragraph{Problem source.}
\#\# Physical scenario

A radar transmitter \textbackslash{}( T \textbackslash{}) is fixed to a reference frame \textbackslash{}( S' \textbackslash{}) that is moving to the right with speed \textbackslash{}( v \textbackslash{}) relative to reference frame \textbackslash{}( S \textbackslash{}). A mechanical timer (essentially a clock) in frame \textbackslash{}( S' \textbackslash{}), having a period \textbackslash{}( 	au\_0 \textbackslash{}) (measured in \textbackslash{}( S' \textbackslash{})), causes transmitter \textbackslash{}( T \textbackslash{}) to emit timed radar pulses, which travel at the speed of light and are received by \textbackslash{}( R \textbackslash{}), a receiver fixed in frame \textbackslash{}( S \textbackslash{}).

\#\# Question

Show that at receiver \textbackslash{}( R \textbackslash{}) the time interval between pulses arriving from \textbackslash{}( T \textbackslash{}) is not \textbackslash{}( 	au \textbackslash{}) or \textbackslash{}( 	au\_0 \textbackslash{}), but \textbackslash{}( 	au\_R = 	au\_0 \textbackslash{}sqrt\{\textbackslash{}frac\{c+v\}\{c-v\}\} \textbackslash{}).

\#\# Answer choices

- A: A: 	au\_0 \textbackslash{}sqrt\{\textbackslash{}frac\{c-v\}\{c-v\}\}

- B: B: 	au\_0 \textbackslash{}left( \textbackslash{}frac\{c+v\}\{c+v\} 
ight)

- C: C: 	au\_0 \textbackslash{}sqrt\{\textbackslash{}frac\{c+v\}\{c-v\}\}

- D: D: \textbackslash{}frac\{	au\_0\}\{\textbackslash{}sqrt\{1 - \textbackslash{}frac\{v\textasciicircum{}2\}\{c\textasciicircum{}2\}\}\}

\#\# Figure caption (auxiliary; use the image as primary evidence)

The image is a diagram representing a physics scenario involving motion and signal transmission. It includes the following components:

1. **Rectangular Frames:** There are two vertical lines labeled \textbackslash{}( S \textbackslash{}) and \textbackslash{}( S' \textbackslash{}). These appear to denote different frames of reference or systems.

2. **Objects within Frames:**
   - The \textbackslash{}( S \textbackslash{}) frame contains a semicircular object labeled \textbackslash{}( R \textbackslash{}).
   - The \textbackslash{}( S' \textbackslash{}) frame contains another semicircular object labeled \textbackslash{}( T \textbackslash{}), and a rectangular block labeled \textbackslash{}( \textbackslash{}tau\_0 \textbackslash{}).

3. **Lines and Arrows:**
   - A horizontal double arrow labeled \textbackslash{}( \textbackslash{}vec\{v\} \textbackslash{}) originates from the \textbackslash{}( S' \textbackslash{}) frame and points to the right, indicating velocity or a direction of motion.

4. **Wave Representation:** A series of curved lines between the \textbackslash{}( S' \textbackslash{}) frame and \textbackslash{}( S \textbackslash{}) frame suggest a wave or signal transmission originating from \textbackslash{}( T \textbackslash{}) towards \textbackslash{}( R \textbackslash{}).

5. **Text and Labels:** The diagram is labeled with the letters \textbackslash{}( A \textbackslash{}) and \textbackslash{}( \textbackslash{}tau\_0 \textbackslash{}), which likely represent specific variables relevant to the depicted scenario.

The diagram illustrates a conceptual model commonly used in physics to discuss relative motion, signal propagation, or interactions between different reference frames.

\#\# Dataset metadata

Relativity; Physical Model Grounding Reasoning, Multi-Formula Reasoning

\paragraph{Recorded answer/context.}
C: C: 	au\_0 \textbackslash{}sqrt\{\textbackslash{}frac\{c+v\}\{c-v\}\}

\paragraph{Figure/image path.}
/root/proposal\_for\_physic/science-mango/phyx\_mini\_run/phyx\_data/test\_image/593.png

\paragraph{Formalization target.}
create a compiling Lean file with sorry bodies at `PhyXMiniProblems/problem\_phyx\_mini\_0593.lean`. The Lean declarations must preserve the physical quantities, dimensions or dimensional roles, figure labels, governing-law hypotheses, and final relation expressed by this problem.
Use Mathlib/Physlib names found through LeanExplore where available. If a physics API is missing, introduce faithful local abstractions rather than scalar placeholder aliases.

\begin{theorem}[Physics formalization target]
\label{thm:physics:phyx_mini_0593:target}
\lean{PhyXMiniProblems.ProblemPhyXMini0593.problem_phyx_mini_0593}
\uses{def:physics:phyx-mini-0593:phyxminiproblems-problemphyxmini0593-durationreadout, def:physics:phyx-mini-0593:phyxminiproblems-problemphyxmini0593-speedreadout, def:physics:phyx-mini-0593:phyxminiproblems-problemphyxmini0593-vacuumlightspeedreadout, def:physics:phyx-mini-0593:phyxminiproblems-problemphyxmini0593-radarpulsetimingsetup, def:physics:phyx-mini-0593:phyxminiproblems-problemphyxmini0593-matchesradarpulsescenario, def:physics:phyx-mini-0593:phyxminiproblems-problemphyxmini0593-matchessuppliedradarfigure, def:physics:phyx-mini-0593:phyxminiproblems-problemphyxmini0593-hasphysicalradarparameters, def:physics:phyx-mini-0593:phyxminiproblems-problemphyxmini0593-satisfiesrelativisticemissionclocklaw, def:physics:phyx-mini-0593:phyxminiproblems-problemphyxmini0593-satisfiesrecedingradarpulsekinematics}
The assigned autoformalize agent should translate this physics problem into Lean declarations in the covered file, with theorem and lemma proof bodies written as `by sorry`.
\end{theorem}
\begin{proof}
This is an autoformalization task, not a proof task. Produce statements that can later be proved by the physics prover without weakening the physical model.
\end{proof}

% --- Archon declaration topology begin ---
\section{Declaration topology}
\begin{definition}[Duration Quantity]
  \label{def:physics:phyx-mini-0593:phyxminiproblems-problemphyxmini0593-durationquantity}
  \lean{PhyXMiniProblems.ProblemPhyXMini0593.DurationQuantity}
  A nonnegative, unit-independent physical duration
\end{definition}
\begin{proof}
  This is the declared model definition of Duration Quantity.
\end{proof}
\begin{definition}[Length Quantity]
  \label{def:physics:phyx-mini-0593:phyxminiproblems-problemphyxmini0593-lengthquantity}
  \lean{PhyXMiniProblems.ProblemPhyXMini0593.LengthQuantity}
  A nonnegative, unit-independent physical length
\end{definition}
\begin{proof}
  This is the declared model definition of Length Quantity.
\end{proof}
\begin{definition}[Speed Quantity]
  \label{def:physics:phyx-mini-0593:phyxminiproblems-problemphyxmini0593-speedquantity}
  \lean{PhyXMiniProblems.ProblemPhyXMini0593.SpeedQuantity}
  A nonnegative, unit-independent physical speed
\end{definition}
\begin{proof}
  This is the declared model definition of Speed Quantity.
\end{proof}
\begin{definition}[duration Readout]
  \label{def:physics:phyx-mini-0593:phyxminiproblems-problemphyxmini0593-durationreadout}
  \lean{PhyXMiniProblems.ProblemPhyXMini0593.durationReadout}
  \uses{def:physics:phyx-mini-0593:phyxminiproblems-problemphyxmini0593-durationquantity}
  Read a physical duration in a selected time unit
\end{definition}
\begin{proof}
  This is the declared model definition of duration Readout.
\end{proof}
\begin{definition}[length Readout]
  \label{def:physics:phyx-mini-0593:phyxminiproblems-problemphyxmini0593-lengthreadout}
  \lean{PhyXMiniProblems.ProblemPhyXMini0593.lengthReadout}
  \uses{def:physics:phyx-mini-0593:phyxminiproblems-problemphyxmini0593-lengthquantity}
  Read a physical length in a selected length unit
\end{definition}
\begin{proof}
  This is the declared model definition of length Readout.
\end{proof}
\begin{definition}[speed Readout]
  \label{def:physics:phyx-mini-0593:phyxminiproblems-problemphyxmini0593-speedreadout}
  \lean{PhyXMiniProblems.ProblemPhyXMini0593.speedReadout}
  \uses{def:physics:phyx-mini-0593:phyxminiproblems-problemphyxmini0593-speedquantity}
  Read a physical speed in coherent selected length and time units
\end{definition}
\begin{proof}
  This is the declared model definition of speed Readout.
\end{proof}
\begin{definition}[vacuum Light Speed Readout]
  \label{def:physics:phyx-mini-0593:phyxminiproblems-problemphyxmini0593-vacuumlightspeedreadout}
  \lean{PhyXMiniProblems.ProblemPhyXMini0593.vacuumLightSpeedReadout}
  Read Physlib's exact vacuum speed of light in coherent selected units
\end{definition}
\begin{proof}
  This is the declared model definition of vacuum Light Speed Readout.
\end{proof}
\begin{definition}[Inertial Frame Label]
  \label{def:physics:phyx-mini-0593:phyxminiproblems-problemphyxmini0593-inertialframelabel}
  \lean{PhyXMiniProblems.ProblemPhyXMini0593.InertialFrameLabel}
  The stationary receiver frame S and receding transmitter frame S'
\end{definition}
\begin{proof}
  This is the declared model definition of Inertial Frame Label.
\end{proof}
\begin{definition}[Radar Apparatus]
  \label{def:physics:phyx-mini-0593:phyxminiproblems-problemphyxmini0593-radarapparatus}
  \lean{PhyXMiniProblems.ProblemPhyXMini0593.RadarApparatus}
  Physical apparatus present in the radar timing experiment
\end{definition}
\begin{proof}
  This is the declared model definition of Radar Apparatus.
\end{proof}
\begin{definition}[Signal Kind]
  \label{def:physics:phyx-mini-0593:phyxminiproblems-problemphyxmini0593-signalkind}
  \lean{PhyXMiniProblems.ProblemPhyXMini0593.SignalKind}
  The electromagnetic signal emitted by the transmitter
\end{definition}
\begin{proof}
  This is the declared model definition of Signal Kind.
\end{proof}
\begin{definition}[Axial Direction]
  \label{def:physics:phyx-mini-0593:phyxminiproblems-problemphyxmini0593-axialdirection}
  \lean{PhyXMiniProblems.ProblemPhyXMini0593.AxialDirection}
  Horizontal directions in the orientation of the supplied figure
\end{definition}
\begin{proof}
  This is the declared model definition of Axial Direction.
\end{proof}
\begin{definition}[Figure Label]
  \label{def:physics:phyx-mini-0593:phyxminiproblems-problemphyxmini0593-figurelabel}
  \lean{PhyXMiniProblems.ProblemPhyXMini0593.FigureLabel}
  Literal text labels visible in the supplied radar schematic
\end{definition}
\begin{proof}
  This is the declared model definition of Figure Label.
\end{proof}
\begin{definition}[Radar Pulse Figure]
  \label{def:physics:phyx-mini-0593:phyxminiproblems-problemphyxmini0593-radarpulsefigure}
  \lean{PhyXMiniProblems.ProblemPhyXMini0593.RadarPulseFigure}
  \uses{def:physics:phyx-mini-0593:phyxminiproblems-problemphyxmini0593-inertialframelabel, def:physics:phyx-mini-0593:phyxminiproblems-problemphyxmini0593-radarapparatus, def:physics:phyx-mini-0593:phyxminiproblems-problemphyxmini0593-axialdirection, def:physics:phyx-mini-0593:phyxminiproblems-problemphyxmini0593-figurelabel}
  Typed transcription of the primary raster. The figure gives qualitative placement and direction, but no numerical speed, distance, or time readout
\end{definition}
\begin{proof}
  This is the declared model definition of Radar Pulse Figure.
\end{proof}
\begin{definition}[Radar Pulse Timing Setup]
  \label{def:physics:phyx-mini-0593:phyxminiproblems-problemphyxmini0593-radarpulsetimingsetup}
  \lean{PhyXMiniProblems.ProblemPhyXMini0593.RadarPulseTimingSetup}
  \uses{def:physics:phyx-mini-0593:phyxminiproblems-problemphyxmini0593-durationquantity, def:physics:phyx-mini-0593:phyxminiproblems-problemphyxmini0593-lengthquantity, def:physics:phyx-mini-0593:phyxminiproblems-problemphyxmini0593-speedquantity, def:physics:phyx-mini-0593:phyxminiproblems-problemphyxmini0593-inertialframelabel, def:physics:phyx-mini-0593:phyxminiproblems-problemphyxmini0593-radarapparatus, def:physics:phyx-mini-0593:phyxminiproblems-problemphyxmini0593-signalkind, def:physics:phyx-mini-0593:phyxminiproblems-problemphyxmini0593-axialdirection, def:physics:phyx-mini-0593:phyxminiproblems-problemphyxmini0593-radarpulsefigure}
  Independent frame assignments and physical observables for two consecutive pulses. In particular, receiverArrivalIntervalTauR is not defined from a Doppler formula or answer choice. The two source--receiver separations and two flight durations retain the light-propagation route by which the initial separation cancels from the requested interval
\end{definition}
\begin{proof}
  This is the declared model definition of Radar Pulse Timing Setup.
\end{proof}
\begin{definition}[speed Fraction Of Light]
  \label{def:physics:phyx-mini-0593:phyxminiproblems-problemphyxmini0593-speedfractionoflight}
  \lean{PhyXMiniProblems.ProblemPhyXMini0593.speedFractionOfLight}
  \uses{def:physics:phyx-mini-0593:phyxminiproblems-problemphyxmini0593-speedreadout, def:physics:phyx-mini-0593:phyxminiproblems-problemphyxmini0593-vacuumlightspeedreadout, def:physics:phyx-mini-0593:phyxminiproblems-problemphyxmini0593-radarpulsetimingsetup}
  The dimensionless relative speed β = v/c, evaluated in SI units
\end{definition}
\begin{proof}
  This is the declared model definition of speed Fraction Of Light.
\end{proof}
\begin{definition}[moving Frame Lorentz Factor]
  \label{def:physics:phyx-mini-0593:phyxminiproblems-problemphyxmini0593-movingframelorentzfactor}
  \lean{PhyXMiniProblems.ProblemPhyXMini0593.movingFrameLorentzFactor}
  \uses{def:physics:phyx-mini-0593:phyxminiproblems-problemphyxmini0593-radarpulsetimingsetup, def:physics:phyx-mini-0593:phyxminiproblems-problemphyxmini0593-speedfractionoflight}
  Physlib's Lorentz factor for the frame speed v/c
\end{definition}
\begin{proof}
  This is the declared model definition of moving Frame Lorentz Factor.
\end{proof}
\begin{definition}[Matches Radar Pulse Scenario]
  \label{def:physics:phyx-mini-0593:phyxminiproblems-problemphyxmini0593-matchesradarpulsescenario}
  \lean{PhyXMiniProblems.ProblemPhyXMini0593.MatchesRadarPulseScenario}
  \uses{def:physics:phyx-mini-0593:phyxminiproblems-problemphyxmini0593-radarpulsetimingsetup}
  Apparatus, frame, triggering, and travel roles stated in the problem
\end{definition}
\begin{proof}
  This is the declared model definition of Matches Radar Pulse Scenario.
\end{proof}
\begin{definition}[Matches Supplied Radar Figure]
  \label{def:physics:phyx-mini-0593:phyxminiproblems-problemphyxmini0593-matchessuppliedradarfigure}
  \lean{PhyXMiniProblems.ProblemPhyXMini0593.MatchesSuppliedRadarFigure}
  \uses{def:physics:phyx-mini-0593:phyxminiproblems-problemphyxmini0593-radarpulsefigure}
  Primary-image evidence: S and R are on the left; S', T, and the timer box labelled τ₀ are on the right; the frame arrow points right while the drawn wavefronts propagate left. The literal baseline label A is recorded without assigning it an unsupported physical interpretation
\end{definition}
\begin{proof}
  This is the declared model definition of Matches Supplied Radar Figure.
\end{proof}
\begin{definition}[Has Physical Radar Parameters]
  \label{def:physics:phyx-mini-0593:phyxminiproblems-problemphyxmini0593-hasphysicalradarparameters}
  \lean{PhyXMiniProblems.ProblemPhyXMini0593.HasPhysicalRadarParameters}
  \uses{def:physics:phyx-mini-0593:phyxminiproblems-problemphyxmini0593-durationreadout, def:physics:phyx-mini-0593:phyxminiproblems-problemphyxmini0593-lengthreadout, def:physics:phyx-mini-0593:phyxminiproblems-problemphyxmini0593-speedreadout, def:physics:phyx-mini-0593:phyxminiproblems-problemphyxmini0593-vacuumlightspeedreadout, def:physics:phyx-mini-0593:phyxminiproblems-problemphyxmini0593-radarpulsetimingsetup}
  Positivity and the ordinary receding, subluminal branch. The strict positive speed is what makes the three intervals distinct; no Doppler factor is stated here
\end{definition}
\begin{proof}
  This is the declared model definition of Has Physical Radar Parameters.
\end{proof}
\begin{definition}[Satisfies Relativistic Emission Clock Law]
  \label{def:physics:phyx-mini-0593:phyxminiproblems-problemphyxmini0593-satisfiesrelativisticemissionclocklaw}
  \lean{PhyXMiniProblems.ProblemPhyXMini0593.SatisfiesRelativisticEmissionClockLaw}
  \uses{def:physics:phyx-mini-0593:phyxminiproblems-problemphyxmini0593-durationreadout, def:physics:phyx-mini-0593:phyxminiproblems-problemphyxmini0593-radarpulsetimingsetup, def:physics:phyx-mini-0593:phyxminiproblems-problemphyxmini0593-movingframelorentzfactor}
  Special-relativistic time dilation between successive ticks of the co-moving timer and their emission events in S. This generic law identifies the intermediate interval τ; it does not state the receiver's interval
\end{definition}
\begin{proof}
  This is the declared model definition of Satisfies Relativistic Emission Clock Law.
\end{proof}
\begin{definition}[Satisfies Receding Radar Pulse Kinematics]
  \label{def:physics:phyx-mini-0593:phyxminiproblems-problemphyxmini0593-satisfiesrecedingradarpulsekinematics}
  \lean{PhyXMiniProblems.ProblemPhyXMini0593.SatisfiesRecedingRadarPulseKinematics}
  \uses{def:physics:phyx-mini-0593:phyxminiproblems-problemphyxmini0593-durationreadout, def:physics:phyx-mini-0593:phyxminiproblems-problemphyxmini0593-lengthreadout, def:physics:phyx-mini-0593:phyxminiproblems-problemphyxmini0593-speedreadout, def:physics:phyx-mini-0593:phyxminiproblems-problemphyxmini0593-vacuumlightspeedreadout, def:physics:phyx-mini-0593:phyxminiproblems-problemphyxmini0593-radarpulsetimingsetup}
  Two-pulse kinematics in the stationary receiver frame. Between emissions the source--receiver separation grows by v τ; each pulse flight obeys distance equals c times flight time; and τ\_R is the second arrival time minus the first. These are governing relations involving independent observables, not the requested square-root Doppler formula
\end{definition}
\begin{proof}
  This is the declared model definition of Satisfies Receding Radar Pulse Kinematics.
\end{proof}
\begin{definition}[Answer Choice]
  \label{def:physics:phyx-mini-0593:phyxminiproblems-problemphyxmini0593-answerchoice}
  \lean{PhyXMiniProblems.ProblemPhyXMini0593.AnswerChoice}
  Labels of the four interval formulas printed in the source
\end{definition}
\begin{proof}
  This is the declared model definition of Answer Choice.
\end{proof}
\begin{definition}[displayed Arrival Interval Factor]
  \label{def:physics:phyx-mini-0593:phyxminiproblems-problemphyxmini0593-displayedarrivalintervalfactor}
  \lean{PhyXMiniProblems.ProblemPhyXMini0593.displayedArrivalIntervalFactor}
  \uses{def:physics:phyx-mini-0593:phyxminiproblems-problemphyxmini0593-speedreadout, def:physics:phyx-mini-0593:phyxminiproblems-problemphyxmini0593-vacuumlightspeedreadout, def:physics:phyx-mini-0593:phyxminiproblems-problemphyxmini0593-radarpulsetimingsetup, def:physics:phyx-mini-0593:phyxminiproblems-problemphyxmini0593-answerchoice}
  The dimensionless factor printed beside each answer label. These definitions record the multiple-choice data; none is assumed to equal the physical arrival interval
\end{definition}
\begin{proof}
  This is the declared model definition of displayed Arrival Interval Factor.
\end{proof}
\begin{definition}[displayed Arrival Interval Readout]
  \label{def:physics:phyx-mini-0593:phyxminiproblems-problemphyxmini0593-displayedarrivalintervalreadout}
  \lean{PhyXMiniProblems.ProblemPhyXMini0593.displayedArrivalIntervalReadout}
  \uses{def:physics:phyx-mini-0593:phyxminiproblems-problemphyxmini0593-durationreadout, def:physics:phyx-mini-0593:phyxminiproblems-problemphyxmini0593-radarpulsetimingsetup, def:physics:phyx-mini-0593:phyxminiproblems-problemphyxmini0593-answerchoice, def:physics:phyx-mini-0593:phyxminiproblems-problemphyxmini0593-displayedarrivalintervalfactor}
  The interval readout represented by a displayed answer choice
\end{definition}
\begin{proof}
  This is the declared model definition of displayed Arrival Interval Readout.
\end{proof}
\begin{definition}[recorded Dataset Answer]
  \label{def:physics:phyx-mini-0593:phyxminiproblems-problemphyxmini0593-recordeddatasetanswer}
  \lean{PhyXMiniProblems.ProblemPhyXMini0593.recordedDatasetAnswer}
  \uses{def:physics:phyx-mini-0593:phyxminiproblems-problemphyxmini0593-answerchoice}
  The answer label recorded by the source dataset
\end{definition}
\begin{proof}
  This is the declared model definition of recorded Dataset Answer.
\end{proof}
\begin{lemma}[receiver Arrival Interval from kinematics]
  \label{lem:physics:phyx-mini-0593:phyxminiproblems-problemphyxmini0593-receiverarrivalinterval-from-kinematics}
  \lean{PhyXMiniProblems.ProblemPhyXMini0593.receiverArrivalInterval_from_kinematics}
  \uses{def:physics:phyx-mini-0593:phyxminiproblems-problemphyxmini0593-durationreadout, def:physics:phyx-mini-0593:phyxminiproblems-problemphyxmini0593-speedreadout, def:physics:phyx-mini-0593:phyxminiproblems-problemphyxmini0593-vacuumlightspeedreadout, def:physics:phyx-mini-0593:phyxminiproblems-problemphyxmini0593-radarpulsetimingsetup, def:physics:phyx-mini-0593:phyxminiproblems-problemphyxmini0593-hasphysicalradarparameters, def:physics:phyx-mini-0593:phyxminiproblems-problemphyxmini0593-satisfiesrecedingradarpulsekinematics}
  The independent two-pulse flight equations first give the receding-source arrival relation τ\_R = τ (c + v) / c. This helper is a consequence of the governing kinematics and contains no proper-time Doppler factor
\end{lemma}
\begin{proof}
  The conclusion follows from the governing relations and figure readouts named in the statement of receiver Arrival Interval from kinematics.
\end{proof}
% --- Archon declaration topology end ---
% --- Archon physics formalization source end ---
